% Part: methods
% Chapter: proofs
% Section: inference-patterns

\documentclass[../../../include/open-logic-section]{subfiles}

\begin{document}

\olfileid{mth}{prf}{pat}

\olsection{నిగమన నమూనాలు}

నిరూపణలు విడివిడి నిగమనాలతో నిర్మితమవుతాయి. నిగమనం చేసినప్పుడు సాధారణంగా ``కాబట్టి'', ``అందువల్ల'', ``అందుచేత'' వంటి పదంతో దాన్ని సూచిస్తాం. ఆ నిగమనం మన నిరూపణలో ఇప్పటికే లభించిన ఒకటి లేదా రెండు విషయాలపై ఆధారపడుతుంది---అది తీసుకున్న పరికల్పన కావచ్చు, ముందరి నిగమనం నుంచి వచ్చిన ఫలితం కావచ్చు. స్పష్టత కోసం వాటికి గుర్తులు పెట్టి, కొత్త నిగమనంలో ఏ ప్రకటనలను ఉపయోగిస్తున్నామో చెప్పవచ్చు. ముందరి విషయాల నుంచి కొత్త నిష్కర్ష \emph{ఎందుకు} వస్తుందో కూడా తరచుగా వివరిస్తాం. నిరూపణలలో చాలా తరచుగా కనిపించే కొన్ని నిగమన నమూనాలను కింద పరిశీలిద్దాం. ఆగమన నిరూపణల వంటి కొన్ని నమూనాలు మరింత విస్తృతమైనవి (వాటిని తరువాత చర్చిస్తాం).

ఒక నిగమన నమూనాను ఇప్పటికే చూశాం: నిర్వచనాన్ని విప్పడం లేదా వర్తింపజేయడం. నిర్వచనాన్ని విప్పినప్పుడు నిర్వచ్యపదం ఉన్న విషయాన్ని నిర్వచక భాగంతో తిరిగి చెబుతాం. ఉదాహరణకు, నిరూపణలో $D = E$ అని ముందే స్థాపించాం అనుకుందాం (a). అప్పుడు సమితుల కోసం $=$ నిర్వచనాన్ని వర్తింపజేసి ఇలా నిగమించవచ్చు: ``అందువల్ల, (a) నుంచి నిర్వచనం ప్రకారం, $D$లోని ప్రతి \tetoken{మూలకం}{!!{element}}, $E$లో \tetoken{ఒక మూలకం}{!!a{element}}; అలాగే వ్యతిరేక దిశలోనూ.''

కొంత గందరగోళంగా అనిపించవచ్చు: నిగమనం చేసే చోట కాక, దానికి ముందే దాని సమర్థనను తరచుగా రాస్తాం. $D = E$ను ఇంకా నిరూపించలేదు, కానీ నిరూపించాలనుకుందాం. $D = E$ మన లక్ష్య నిష్కర్ష అయితే, నిర్వచనాన్ని వర్తింపజేసి లక్ష్యాన్ని మరోలా చెప్పవచ్చు: $D = E$ను నిరూపించాలంటే $D$లోని ప్రతి \tetoken{మూలకం}{!!{element}}, $E$లో \tetoken{ఒక మూలకం}{!!a{element}} అని, వ్యతిరేక దిశలోనూ అలాగే అని నిరూపించాలి. కాబట్టి నిరూపణకు ఈ రూపం ఉంటుంది: (a) $D$లోని ప్రతి \tetoken{మూలకం}{!!{element}}, $E$లో \tetoken{ఒక మూలకం}{!!a{element}} అని నిరూపించడం; (b) $E$లోని ప్రతి \tetoken{మూలకం}{!!{element}}, $D$లో \tetoken{ఒక మూలకం}{!!a{element}} అని నిరూపించడం; (c) అందువల్ల (a), (b) నుంచి $=$ నిర్వచనం ప్రకారం $D = E$. అయితే సాధారణంగా మనం ఇలా రాయకుండా ఈ విధంగా రాస్తాం:
\begin{quote}
$D = E$ అని చూపాలి. $=$ నిర్వచనం ప్రకారం, ఇది $D$లోని ప్రతి \tetoken{మూలకం}{!!{element}}, $E$లో \tetoken{ఒక మూలకం}{!!a{element}} అని, అలాగే వ్యతిరేక దిశలోనూ చూపడమే.

(a) \dots ($D$లోని ప్రతి \tetoken{మూలకం}{!!{element}}, $E$లో \tetoken{ఒక మూలకం}{!!a{element}} అని నిరూపణ) \dots

(b) \dots ($E$లోని ప్రతి \tetoken{మూలకం}{!!{element}}, $D$లో \tetoken{ఒక మూలకం}{!!a{element}} అని నిరూపణ) \dots
\end{quote}

\subsection{సంయోగాన్ని ఉపయోగించడం}

బహుశా అతి సరళమైన నిగమన నమూనా, సంయోగంలోని ఒక సంయోజ్యాన్ని నిష్కర్షగా తీసుకోవడం. మరో మాటలో, $p$ మరియు~$q$ అని పరికల్పనగా తీసుకున్నా లేదా ముందే నిరూపించినా, $p$ను (అలాగే $q$ను) నిగమించవచ్చు. ఇది ఎంత ప్రాథమిక నిగమనమంటే తరచుగా ప్రత్యేకంగా చెప్పరు. ఉదాహరణకు $D = E$ నిర్వచనాన్ని విప్పిన తరువాత $D$లోని ప్రతి \tetoken{మూలకం}{!!{element}}, $E$లో \tetoken{ఒక మూలకం}{!!a{element}} అని, అలాగే వ్యతిరేక దిశలోనూ అని స్థాపించాం. దాని నుంచి $E$లోని ప్రతి \tetoken{మూలకం}{!!{element}}, $D$లో \tetoken{ఒక మూలకం}{!!a{element}} అని నిగమించవచ్చు (అదే ``వ్యతిరేక దిశలోనూ'' అన్న భాగం).

\subsection{సంయోగాన్ని నిరూపించడం}

కొన్నిసార్లు నిరూపించాల్సింది సంయోగ రూపంలో ఉంటుంది: ``$p$ మరియు $q$ను నిరూపించండి'' అని అడుగుతారు. అప్పుడు రెండు పనులే చేయాలి: $p$ను నిరూపించి, తరువాత $q$ను నిరూపించాలి. స్పష్టత కోసం నిరూపణను రెండు భాగాలుగా విడదీసి వాటికి గుర్తులు పెట్టవచ్చు. మొదటి గమనికల్లో కాగితం పైభాగంలో ``(1) $p$ను నిరూపించాలి'', మధ్యలో ``(2) $q$ను నిరూపించాలి'' అని రాయవచ్చు. (సంయోగాన్ని నిరూపించమని నేరుగా అడగకపోయినా, నిరూపణకు అది అవసరమవవచ్చు. ఉదాహరణకు $D = E$ను నిరూపించమంటే, $=$ నిర్వచనాన్ని విప్పిన తరువాత $D$లోని ప్రతి \tetoken{మూలకం}{!!{element}}, $E$లో \tetoken{ఒక మూలకం}{!!a{element}} అని \emph{మరియు} $E$లోని ప్రతి \tetoken{మూలకం}{!!{element}}, $D$లో \tetoken{ఒక మూలకం}{!!a{element}} అని నిరూపించాల్సి ఉంటుందని తెలుస్తుంది.)

\subsection{వికల్పాన్ని నిరూపించడం}

నిరూపించాల్సింది వికల్ప రూపంలో (అంటే ``$p$ లేదా $q$'' రూపంలోని ప్రకటనగా) ఉంటే, దాని వికల్పాలలో ఒకటి సత్యమని చూపితే చాలు. అయితే సిద్ధాంతపు పరికల్పనల నుంచే ఏదో ఒక వికల్పం నేరుగా రావడం సాధారణంగా జరగదు. ఎక్కువగా సిద్ధాంతపు పరికల్పనలే వికల్ప రూపంలో ఉంటాయి; లేదా ఒక రకానికి చెందిన అన్ని వస్తువులకు రెండు లక్షణాలలో ఒకటి ఉంటుందని చూపుతుంటారు, కాని కొన్ని వస్తువులకు మొదటి లక్షణం, మరికొన్నింటికి రెండోది ఉంటుంది. అలాంటప్పుడు సందర్భాలవారీ నిరూపణ ఉపయోగపడుతుంది (కింద చూడండి).

\subsection{సోపాధిక నిరూపణ}

మీకు ఎదురయ్యే అనేక సిద్ధాంతాలు సోపాధిక రూపంలో ఉంటాయి (అంటే $p$ ఉంటే $q$ కూడా సత్యమని చూపాలి). వాటిని అమర్చడం సులభం---సోపాధిక పూర్వపక్షాన్ని (ఇక్కడ $p$ను) పరికల్పనగా తీసుకొని దాని నుంచి నిష్కర్ష~$q$ను నిరూపించండి. మీ సిద్ధాంతం ``$p$ అయితే $q$'' అనుకుంటే, ``$p$ అనుకుందాం'' అని నిరూపణను మొదలుపెట్టి, చివరికి~$q$ను నిరూపించి ఉండాలి.

సోపాధికాలను వేర్వేరు మాటల్లో చెప్పవచ్చు. ``$p$ అయితే $q$''కు బదులు ``$p$ నిజమవాలంటే $q$ తప్పనిసరి'', ``$q$, ఒకవేళ $p$ అయితే'', లేదా ``$p$ అన్న షరతుపై $q$'' అని సిద్ధాంతం చెప్పవచ్చు. ఇవన్నీ ఒకే అర్థం ఇస్తాయి: $p$ను అనుకొని, దాని నుంచి~$q$ను నిరూపించాలి. ద్విసోపాధికం (``$p$ అప్పుడూ అప్పుడే $q$'') నిజానికి రెండు సోపాధికాల కలయిక: $p$ అయితే $q$, అలాగే $q$ అయితే~$p$. కనుక మొదటి సోపాధికానికి ఒకసారి, రెండవ దానికి మరొకసారి సోపాధిక నిరూపణ చేయాలి. అయితే కొన్నిసార్లు ``అప్పుడూ అప్పుడే'' ప్రకటనలను గొలుసుగా కలిపి ``$p$'' నుంచి మొదలై ``$q$''తో ముగించి ద్విసోపాధికాన్ని నిరూపించవచ్చు---అప్పుడు ప్రతి దశ నిజంగానే ``అప్పుడూ అప్పుడే'' అని నిర్ధారించాలి.

\subsection{సార్వత్రిక వాదనలు}

సార్వత్రిక వాదనను ఉపయోగించడం సులభం: ఏ వస్తువుకైనా ఒక విషయం సత్యమైతే ప్రతి నిర్దిష్ట వస్తువుకూ అది సత్యమే. ఉదాహరణకు మీ నిరూపణ పరికల్పన $A \subseteq B$ అయితే, $\subseteq$ నిర్వచనాన్ని విప్పితే ప్రతి $x \in A$కు $x \in B$ అని అర్థం. కాబట్టి $z \in A$ అని ముందే తెలిస్తే, $z \in B$ అని నిగమించవచ్చు.

సార్వత్రిక వాదనను నిరూపించడం కాస్త కష్టంగా అనిపించవచ్చు. సాధారణంగా అవి ``$x$కు~$P$ ఉంటే, దానికి~$Q$ ఉంటుంది'' లేదా ``అన్ని $P$s, $Q$s'' అనే రూపంలో ఉంటాయి. ఈ రూపానికి అచ్చంగా సరిపోకపోవచ్చు; నిరూపించమన్నది ఖచ్చితంగా ఏమిటో గ్రహించడానికి కొంత సాధన కావాలి. కానీ తరచుగా ఒక లక్షణం గల అన్ని వస్తువులకు మరో లక్షణమూ ఉందని నిరూపించాలి.

సార్వత్రిక వాదనను నిరూపించడానికి మొదటి లక్షణం గల వస్తువులకు పేర్లు లేదా చరరాశులను పరిచయం చేసి, వాటికి రెండవ లక్షణమూ ఉందని చూపాలి. \emph{అన్ని}~$P$s గురించి నిరూపించాలంటే \emph{ఏదైనా ఒక}~$P$ గురించి నిరూపించాలి అని చెప్పవచ్చు. ప్రవేశపెట్టిన పేరు, అటువంటి ఏదైనా~$P$కు పేరు. సాధారణంగా ఇలాంటి వస్తువులకు ఒకే అక్షరాన్ని వాడుతాం; అక్షరాల వాడుకలో సంప్రదాయాలూ ఉన్నాయి: ఉదా., సహజ సంఖ్యలకు $n$, \tetoken{సూత్రాలకు}{!!{formula}s} $!A$, సమితులకు $A$, ప్రమేయాలకు $f$ మొదలైనవి.

ఇందులో చాతుర్యం నిరూపణ అంతటా సాధారణత్వాన్ని నిలపడం. ఏదైనా ఒక వస్తువుకు (``$x$''కు)~$P$ లక్షణం ఉందని తీసుకొని, నిర్వచనాలు లేదా అనుమతించిన పరికల్పనల ఆధారంగానే $x$కు~$Q$ లక్షణమూ ఉందని చూపాలి. $P$ లక్షణం ఉండటాన్ని తప్పించి $x$ ఏమిటో ప్రత్యేకంగా నిర్ణయించలేదు కాబట్టి, $P$ ఉన్న ప్రతిదానికీ~$Q$ ఉంటుందని చెప్పవచ్చు. సంక్షిప్తంగా, $x$ అనేది~$P$ లక్షణం గల \emph{అన్ని} వస్తువులకు ప్రతినిధి.

\begin{prop}
  అన్ని సమితులు $A$, $B$లకు $A \subseteq A \cup B$.
\end{prop}

\begin{proof}
  $A$, $B$ ఏవైనా సమితులనుకుందాం. $A
  \subseteq A \cup B$ అని చూపాలి. $\subseteq$ నిర్వచనం ప్రకారం, ప్రతి $x$కు $x \in A$ అయితే $x \in A \cup B$ అని చూపడమే. కాబట్టి $A$లోని ఏదైనా \tetoken{మూలకం}{!!{element}}గా $x \in A$ను తీసుకుందాం. $x \in A
  \cup B$ అని చూపాలి. $x \in A$ కనుక, $x \in A$ లేదా $x \in B$. అందువల్ల $x \in
  \Setabs{x}{x \in A \lor x \in B}$. అది $\cup
  $ నిర్వచనం ప్రకారం $x \in A \cup B$ అని అర్థం.
\end{proof}

\subsection{సందర్భాలవారీ నిరూపణ}

పరికల్పనగా లేదా ముందే స్థాపించిన నిష్కర్షగా ఒక వికల్పం ఉందనుకుందాం---$p$ లేదా $q$ సత్యమని మీరు అనుకున్నారు లేదా నిరూపించారు. ఇప్పుడు $r$ను నిరూపించాలి. రెండు దశల్లో చేయండి: మొదట $p$ సత్యమని తీసుకొని~$r$ను నిరూపించండి; తరువాత $q$ సత్యమని తీసుకొని మళ్లీ~$r$ను నిరూపించండి. రెండు ప్రత్యామ్నాయాలలో ఒకటి సత్యమని మనకు తెలుసు లేదా పరికల్పనగా ఉంది కాబట్టి ఇది పనిచేస్తుంది. ఏ ప్రత్యామ్నాయమైనా~$r$ సత్యానికి సరిపోతుందని ఈ రెండు దశలు చూపుతాయి. (రెండూ సత్యమైతే $r$~ఎందుకు సత్యమో చెప్పడానికి ఒకటి కాక రెండు కారణాలు ఉన్నట్లే. $p$, $q$ రెండింటినీ కలిపి పరికల్పనగా తీసుకొని $r$~సత్యమని విడిగా నిరూపించాల్సిన అవసరం లేదు.) మనం చేసేది సూచించడానికి ``సందర్భాలను వేరుచూద్దాం'' అని ప్రకటిస్తాం. ఉదాహరణకు $x \in B \cup C$ అని తెలుసనుకుందాం. $B \cup C$ నిర్వచనం $\Setabs{x}{x \in B \text{ or } x \in C}$. మరో మాటలో, నిర్వచనం ప్రకారం $x \in B$ లేదా $x \in C$. దీని నుంచి $x \in A$ను నిరూపించాలంటే మొదట $x \in B$ అని తీసుకొని దాని నుంచి $x \in A$ను నిరూపిస్తాం; తరువాత $x \in C$ అని తీసుకొని మళ్లీ $x \in A$ను నిరూపిస్తాం. పరికల్పన కింద ``సందర్భాలను వేరుచూద్దాం'' అని, దాని కింద ``సందర్భం (1): $x \in B$'', పేజీ మధ్యలో ``సందర్భం (2): $x \in C$'' అని రాయవచ్చు. తరువాత పేజీ పై, కింది భాగాల్లో ఆయా నిరూపణలను నింపాలి.

నిరూపించాల్సిందే వికల్ప రూపంలో ఉన్నప్పుడు సందర్భాలవారీ నిరూపణ ప్రత్యేకంగా ఉపయోగపడుతుంది. ఒక సరళ ఉదాహరణ:

\begin{prop}
$B \subseteq D$, $C \subseteq E$ అనుకుందాం. అప్పుడు $B \cup C \subseteq
D \cup E$.
\end{prop}

\begin{proof}
  (a) $B \subseteq D$, (b) $C \subseteq E$ అని అనుకుందాం. నిర్వచనం ప్రకారం ప్రతి $x \in B$ కూడా $\in D$ (c); ప్రతి $x \in C$ కూడా $\in E$ (d). $B \cup C \subseteq D \cup E$ అని చూపాలంటే, $x \in B \cup C$ అయితే $x \in D \cup E$ అని చూపాలి ($\subseteq$ నిర్వచనం). $x \in B \cup C$ అప్పుడూ అప్పుడే $x \in B$ లేదా $x \in
  C$ ($\cup$ నిర్వచనం). అలాగే $x \in D \cup E$ అప్పుడూ అప్పుడే $x \in
  D$ లేదా $x \in E$. కాబట్టి ప్రతి $x$కు $x \in B$ లేదా $x \in C$ అయితే $x \in D$ లేదా $x \in E$ అని చూపాలి.

  \begin{quote}
  ఇప్పటివరకు మనం నిర్వచనాలను మాత్రమే విప్పాం!{} $\subseteq$, $\cup$ లేకుండా ప్రతిపాదనను తిరిగి చెప్పాం; ఇప్పుడు సార్వత్రిక సోపాధిక వాదనను నిరూపించాలి. పైన చెప్పినట్లు, $x$ అనేది ``అయితే'' భాగం సత్యమని తీసుకునే వస్తువుగా భావించి, ``అప్పుడు'' భాగమూ సత్యమని చూపాలి. మరోలా చెప్పాలంటే $x
  \in B$ లేదా $x \in C$ అని తీసుకొని, $x \in D$ లేదా $x \in
  E$ అని చూపాలి.\footnote{ఈ పేరా మనం చేస్తున్న పనిని వివరిస్తుంది మాత్రమే---నిరూపణలో భాగం కాదు; మీ నిరూపణలు రాసేటప్పుడు ఇంత వివరంగా చెప్పాల్సిన అవసరం లేదు.}
  \end{quote}

  $x \in B$ లేదా $x \in C$ అనుకుందాం. $x \in D$ లేదా $x \in E$ అని చూపాలి. సందర్భాలను వేరుచూద్దాం.

  సందర్భం 1: $x \in B$. (c) ప్రకారం $x \in D$. అందువల్ల $x \in D$ లేదా $x \in
  E$. (ఇక్కడ ముందరి ఉపవిభాగంలో చర్చించిన నిగమనం చేశాం!)

  సందర్భం 2: $x \in C$. (d) ప్రకారం $x \in E$. అందువల్ల $x \in D$ లేదా $x \in E$.
\end{proof}

\subsection{అస్తిత్వ వాదనను నిరూపించడం}

అస్తిత్వ వాదనను నిరూపించమంటే ప్రశ్న సాధారణంగా ``ఏదో~$x$కు $\dots x \dots$ అని నిరూపించండి'' అనే రూపంలో ఉంటుంది; అంటే ``$\dots x \dots$'' సూచించే లక్షణం గల ఏదో వస్తువు ఉందని చూపాలి. అప్పుడు తగిన వస్తువును గుర్తించి, దానికి కావలసిన లక్షణం ఉందని నిరూపించాలి. ఇది సరళంగా వినిపించినా, ఇలాంటి నిరూపణ కష్టమవచ్చు. సాధారణంగా ఒక వస్తువును \emph{నిర్మించడం} లేదా \emph{నిర్వచించడం}, తరువాత అలా నిర్వచించిన వస్తువుకు కావలసిన లక్షణం ఉందని నిరూపించడం అవసరం. సరైన వస్తువును కనుక్కోవడం కష్టమవచ్చు; దాని లక్షణాన్ని నిరూపించడం కష్టమవచ్చు; కొన్నిసార్లు వస్తువునే నిజంగా నిర్వచించగలిగామా అని చూపడమూ కష్టమే!{}

సాధారణంగా వస్తువును స్పష్టంగా పేర్కొంటారు, ఉదా., ``$x$ను \dots గా తీసుకుందాం'' (ఇక్కడ \dots{} ఏ వస్తువో చెబుతుంది); అవసరమైతే $\dots$ నిజంగా ఉన్న వస్తువును వివరిస్తుందని నిరూపించి, తరువాత $x$కు~$Q$ లక్షణం ఉందని చూపుతారు. ఒక సరళ ఉదాహరణ ఇదిగో.

\begin{prop}
  $x \in B$ అనుకుందాం. అప్పుడు $A \subseteq
  B$, $A \neq \emptyset$ అయ్యే~$A$ ఒకటి ఉంటుంది.
\end{prop}

\begin{proof}
  $x \in B$ అనుకుందాం. $A = \{x\}$గా తీసుకుందాం.
  \begin{quote}
    ఇక్కడ $A$ సమితిని దాని \tetoken{మూలకాలను}{!!{element}s} జాబితా చేసి నిర్వచించాం. $x$ ఒక వస్తువు అని తీసుకున్నాం; దాని \tetoken{మూలకాలను}{!!{element}s} జాబితా చేసి ఎప్పుడూ ఒక సమితిని ఏర్పరచవచ్చు. కాబట్టి ఇక్కడ $A$ సమితిని నిజంగా నిర్వచించామని విడిగా చూపాల్సిన అవసరం లేదు. అయితే $A$కు ప్రతిపాదన కోరిన లక్షణాలు ఉన్నాయని ఇంకా చూపాలి. అది లేకపోతే నిరూపణ పూర్తి కాదు!{}
  \end{quote}
  $x \in A$ కనుక $A \neq \emptyset$.
  \begin{quote}
    ఇది $A$ను $\{x\}$గా నిర్వచించడంపైనా, $x \in \{x\}$ మరియు $x \notin \emptyset$ అనే స్పష్టమైన విషయాలపైనా ఆధారపడుతుంది.
  \end{quote}
  $\{x\}$లో $x$ మాత్రమే \tetoken{మూలకం}{!!{element}}; ఇంకా $x \in B$. కాబట్టి $A$లోని ప్రతి \tetoken{మూలకం}{!!{element}}, $B$లో \tetoken{ఒక మూలకం}{!!a{element}}. $\subseteq$ నిర్వచనం ప్రకారం $A \subseteq B$.
\end{proof}

\subsection{అస్తిత్వ వాదనలను ఉపయోగించడం}

ఏదో అస్తిత్వ వాదన సత్యమని మీకు తెలుసనుకుందాం (దాన్ని నిరూపించారు లేదా దాన్ని పరికల్పనగా ఉపయోగించవచ్చు); ఉదా., ``ఏదో~$x$కు $x \in A$'' లేదా ``$x \in A$ అయ్యే వస్తువు ఉంది.'' నిరూపణలో దాన్ని ఉపయోగించాలంటే, అది ఉన్నదని పరికల్పన చెప్పే వస్తువులలో ఒకదానికి పేరు ఉన్నట్లు భావించవచ్చు. $A$లో కనీసం ఒక వస్తువు ఉంది కాబట్టి ఆ పేరు సూచించగల వస్తువులు ఉన్నాయి. దేనిని ఎంచుకున్నారో ప్రత్యేకంగా గుర్తించలేకపోవచ్చు, అది $\in A$ అని తప్ప మరేమీ వివరించలేకపోవచ్చు. అయినా నిరూపణ కోసం దాన్ని ఎంచుకొని పేరు పెట్టినట్లు భావించవచ్చు. ఇంతకుముందు వాడని (పరికల్పనలలోనూ లేని) పేరును ఎంచుకోవడం ముఖ్యం; లేకపోతే తప్పులు జరగవచ్చు. నిరూపణలో ``ఏదో $x$కు $x \in A$'' నుంచి ``$a
\in A$ అనుకుందాం'' అని సాగుతాం. తరువాత~$a$ గురించి, ఇతర పరికల్పనలతో కలిపి తర్కించి, ఒక నిష్కర్ష $p$కు చేరుతాం. $p$లో ఇక~$a$ ప్రస్తావన లేకపోతే, $p$ అనేది $a \in A$ అన్న ప్రత్యేక ఎంపికపై ఆధారపడదు; అది ``ఏదో $x$కు $x \in A$'' అనే పరికల్పన నుంచే వస్తుందని చూపినట్లవుతుంది.

\begin{prop}
$A \neq \emptyset$ అయితే, $A \cup B \neq \emptyset$.
\end{prop}

\begin{proof}
  $A \neq \emptyset$ అనుకుందాం. కాబట్టి ఏదో $x$కు $x \in A$.
  \begin{quote}
    ఇక్కడ మొదట ప్రతిపాదన పరికల్పనను మరోలా చెప్పాం. ఆ పరికల్పన, అంటే $A \neq \emptyset$, అస్తిత్వ వాదనను దాచుకుంది; కొన్ని నిర్వచనాలను విప్పినప్పుడే అది తెలుస్తుంది. $=$ నిర్వచనం ప్రకారం $A = \emptyset$ అప్పుడూ అప్పుడే ప్రతి $x \in A$ కూడా $\in \emptyset$, మరియు ప్రతి $x \in
    \emptyset$ కూడా $\in A$. రెండు వైపులనూ నిషేధిస్తే: $A \neq
    \emptyset$ అప్పుడూ అప్పుడే ఏదో $x \in A$కు $\notin \emptyset$, లేదా ఏదో $x \in \emptyset$కు $\notin A$. $\in \emptyset$ అయ్యేది ఏదీ లేదు కనుక రెండవ వికల్పం సత్యం కాదు; ``$x \in A$ మరియు $x
    \notin \emptyset$'' అనేది కేవలం $x \in A$గా సరళమవుతుంది. అందుచేత $A \neq
    \emptyset$ అప్పుడూ అప్పుడే ఏదో $x$కు $x \in A$. ఇది అస్తిత్వ వాదన. ఇప్పుడు $A$లోని \tetoken{మూలకాలలో}{!!{element}s} ఒకదానికి పేరు పెట్టి ఆ వాదనను వాడతాం:
    \sourcecorrection{OLTEMTHPRFPAT-001}{మూల వివరణ చివర $x\neq\emptyset$ అని చెప్పింది; ఇక్కడ నిరూపిస్తున్న సమానార్థకత $A\neq\emptyset$ గురించే. సమితి $A$ను సూచించే సరైన వ్యక్తీకరణతో సరిచేశాం.}
  \end{quote}
  $a \in A$గా తీసుకుందాం.
  \begin{quote}
    ఇప్పుడు $\in A$ అయ్యే వస్తువులలో ఒకదానికి పేరు పెట్టాం. ఇక~$a$ గురించి తర్కిస్తాం; అయితే $a \in A$ అని తప్ప మరేదీ అనుకోకుండా జాగ్రత్తపడతాం:
  \end{quote}
  $a \in A$ కనుక $\cup$ నిర్వచనం ప్రకారం $a \in A \cup B$. అందువల్ల ఏదో $x$కు $x \in A \cup B$; అంటే $A \cup B \neq \emptyset$.
  \begin{quote}
    చివరి దశలో ``$a \in A \cup B$'' నుంచి ``ఏదో $x$కు $x \in A \cup B$''కు వెళ్లాం. రెండవ వాక్యంలో $a$ లేదు; కాబట్టి ``ఏదో $x$కు $x \in A \cup B$'' అనేది ``ఏదో $x$కు $x \in A$'' నుంచే వస్తుందని తెలుసు. అంటే $A \cup B \neq
    \emptyset$.
  \end{quote}
\end{proof}

``$x$'' వంటి బద్ధ చరరాశులను $a$ వంటి ఊహా పేర్ల నుంచి వేరుగా ఉంచడం మంచి అలవాటు కావచ్చు; పైన అలాగే చేశాం. కానీ ఆచరణలో తరచుగా ఒకే $x$ను ఇలా వాడుతాం:
\begin{quote}
$A \neq \emptyset$ అనుకుందాం; అంటే ఏదో $x \in A$ ఉంది. $\cup$ నిర్వచనం ప్రకారం $x \in A \cup B$. కాబట్టి $A \cup B \neq \emptyset$.
\end{quote}
అలా చేసినప్పుడు, వేర్వేరు అస్తిత్వ వాదనలకు వేర్వేరు $x$లు, $y$లు వాడుతున్నామని ప్రత్యేక జాగ్రత్త అవసరం. ఉదాహరణకు, ఈ క్రింది నిరూపణ ``$A \neq
\emptyset$, $B \neq \emptyset$ అయితే $A \cap B \neq \emptyset$'' అనే అసత్య వాదనకు \emph{సరైనది కాదు}.
\begin{quote}
$A \neq \emptyset$, $B \neq \emptyset$ అనుకుందాం. కాబట్టి ఏదో $x$కు $x
\in A$, అలాగే ఏదో $x$కు $x \in B$. $x \in A$ మరియు $x \in
B$ కాబట్టి $\cap$ నిర్వచనం ప్రకారం $x \in A \cap B$. అందువల్ల $A \cap B \neq
\emptyset$.
\end{quote}
తప్పు దశ ఎక్కడ ఉందో గుర్తించి, ఆ ఫలితం ఎందుకు రాదో వివరించగలరా?

\end{document}
